\begin{abstract}
Deploying land-cover classification systems across geographically diverse regions demands models that incrementally learn new classes from new domains without forgetting and with minimal additional storage overhead.
We study this \emph{cross-domain multimodal class-incremental learning} (CIL) problem on HSI+LiDAR data spanning three geographic domains (MUUFL, Trento, Houston~2013), 9 sequential tasks, and 32 classes.
Through systematic drift analysis with six architectures of varying branch coupling, we reveal a previously unreported \emph{architecture-dependent drift asymmetry}: in independent-branch fusion backbones, spatial features drift up to 5.8$\times$ more than spectral features across domains, whereas this asymmetry vanishes in tightly coupled architectures.
We identify three correlated factors---branch independence, parameter capacity asymmetry, and feature coupling---and argue that this asymmetry constitutes a \emph{structural prior} of independent-branch designs that can be exploited as a design signal for non-uniform adaptation.
Based on this insight, we propose \methodname{} (Cross-Modal Compensated Drift LoRA), which adopts a tri-branch Mamba backbone with a separate spectral stream and Siamese spatial processing to preserve the exploitable asymmetry, freezes the spectral branch as a stable anchor after warm-up, and applies asymmetric low-rank adapters to the two spatial branches with rank matched to each branch's effective adaptation dimensionality.
A per-domain whitening normalization (SHINE) and drift-compensated prototype reconstruction (DCR) further address cross-domain feature misalignment.
On the 9-task benchmark (32 classes across 3 domains: 11+6+15), \methodname{} achieves $84.5\% \pm 0.9$ balanced accuracy \emph{without maintaining a replay buffer}, storing only compact prototypes, adapter weights, and normalization statistics (${\sim}75\times$ less extra storage than replay buffers) while leveraging within-domain training data for adapter fine-tuning and prototype recomputation---a natural assumption in fixed-sensor remote sensing deployments. See Section~V for the full data access protocol.
It surpasses all replay-free baselines by over $40$\,pp and is competitive with the best replay methods (iCaRL $77.9\%$, FOSTER $77.0\%$ on their original backbones).
Evaluation across all six domain orderings confirms consistent advantages over replay baselines ($80.9\%$ vs.\ $77.8\%$ and $77.2\%$, winning 5 of 6 orderings).
\end{abstract}
